\chapter{Optimization and memory-efficient training}
A memory optimization is correct only relative to a declared training procedure.
This chapter treats the procedure as an invariant: fix the objective, parameters,
randomness policy, and optimizer-update clock, then change how the computation is
scheduled. We derive weighted accumulation, compare every gradient, verify Adam's
moments and updated parameters, and deliberately break stochastic recomputation.

The executable experiment is a small CPU Float64 MLP, not DOGMA or Hermon DNA
training, a mixed-precision certification, or a hardware benchmark. Its value is a
sharp correctness oracle that later implementations must match where their
assumptions apply. Chapter 26's distinction between equal outputs and equal
gradient graphs remains essential.

\section{Persistent arrays are not peak memory}
\Plate{01-memory}{Memory has multiple owners and lifetimes. Static parameter and
optimizer inventories do not include saved activations, temporary workspaces, or
allocator reservation.}{evod27:memory}
For $P$ parameters, parameter/gradient widths $b_\theta,b_g$, two Adam moments of
width $b_m$, and optional master-copy width $b_{\rm master}$,
\[
M_{\rm arrays}=P(b_\theta+b_g+2b_m+b_{\rm master}).
\]
The default all-FP32 layout gives $16P$ bytes. The reference MLP has
$8\cdot16+16+16\cdot3+3=195$ parameters: this hypothetical inventory is 3,120 bytes.
Its actual Float64 parameter, gradient, and moment arrays would total 6,240 bytes
after all arrays exist, before object overhead or temporary tensors.

A low-precision layout with two-byte weights, two-byte gradients, four-byte
moments, and four-byte master weights also totals $16P$. Lowering the weight dtype
alone does not prove lower total training memory. Peak memory additionally depends
on activation lifetimes, kernel workspaces, optimizer temporaries, cached allocations,
and communication. These components need measurement, not a formula substitution.

\section{Define the reduction before choosing microbatches}
Let $p_{ij}(\theta)$ predict output $j$ for record $i$, and let fixed
nonnegative weights $w_{ij}$ include both validity masks and importance weights.
Define
\[
Z=\sum_{i,j}w_{ij}>0,\qquad
L(\theta)=\frac1Z\sum_{i,j}w_{ij}(p_{ij}(\theta)-y_{ij})^2.
\]
This \Term{weighted reduction} treats a mask as a zero weight, not a target value. Target, prediction, and weight shapes
must match exactly: silent broadcasting can turn a three-output objective into a
different computation while still producing a scalar.

\Listing{weighted_sum}
The batch validator requires finite floating inputs with matching dtype/device,
fixed weights without gradients, no negative weights, and a finite positive
global mass $Z$. Individual finite weights can sum to infinity, so the sum itself
must be checked. Intermediate predictions or squared errors can still overflow;
finite input validation is not a universal numerical-stability guarantee.

For cross-entropy over genomic tokens, replace squared error with token loss
and specify whether the denominator counts valid tokens, sequences, or importance
mass. Averaging each sequence mean weights short and long sequences equally;
averaging all valid token losses does not. Neither is inherently wrong, but they
are different experiments.

\section{Accumulate numerators with one global denominator}
Partition records into disjoint microbatches $B_k$. Keep $\theta$ fixed throughout
the accumulation window and define
\[
N_k=\sum_{i\in B_k,j}w_{ij}(p_{ij}-y_{ij})^2,\qquad
L=\sum_k N_k/Z.
\]
Linearity of differentiation gives
\[
\nabla_\theta L=\sum_k\nabla_\theta(N_k/Z).
\]
This \Term{gradient accumulation} adds contributions in parameter gradient buffers.
A microbatch with zero valid weight contributes zero; it does not divide by a local
zero denominator. The complete window must still have positive mass.

\Plate{02-accum}{For a 5/5/2 record split, valid weight masses are 12, 13, and 6.
Every numerator divides by 31, not by its local mass and then by three.}{evod27:accum}
\Listing{accumulated_grad}
The fixed fixture has 12 records, eight features, and three outputs.
Some outputs are masked and selected first outputs have weight two; total mass
is 31. The batch loss is $1.099552069480512$.

\begin{center}\input{\Generated/microbatches.tex}\end{center}
Microbatch size five creates a partial last microbatch. The measured errors are
maximum absolute differences over \emph{all} parameter gradients. Float64 roundoff
changes addition order, so a tolerance is appropriate; bitwise equality is not
the mathematical contract.

The shortcut $\frac1K\sum_k N_k/Z_k$ changes the objective unless local masses
$Z_k$ are equal or special cancellations occur. Its gradient differs from the
correct gradient by about $0.141178$ in this fixture. This is an executable
counterexample, not an argument that every equal-sized record split is wrong:
equal record counts can still have unequal valid token or weight mass.

\section{An independent derivative oracle}
For this MLP, $H=\tanh(XW_1^\top+b_1)$ and $P=HW_2^\top+b_2$. Holding $H$ fixed
for the last-layer derivative,
\[
D=\frac{2}{Z}W\odot(P-Y),\qquad
\nabla_{W_2}L=D^\top H,\qquad
\nabla_{b_2}L=\sum_iD_{i,:}.
\]
Tests compare these manual expressions with autograd without calling the
accumulation helper. A separate scalar loop computes the loss element by element.
A central finite difference checks one first-layer weight, exercising the nonlinear
chain rule through $\tanh$.

Scaling every weight by the same positive constant leaves $L$ and its gradient
unchanged because numerator and denominator scale together. This metamorphic test
also checks the zero-valid-microbatch case. It is stronger than comparing two
implementations that accidentally share the same reduction mistake.

\section{The optimizer clock advances once per window}
For Adam without weight decay or AMSGrad,
\[
m_t=\beta_1m_{t-1}+(1-\beta_1)g_t,\quad
v_t=\beta_2v_{t-1}+(1-\beta_2)g_t^2,
\]
\[
\widehat m_t=\frac{m_t}{1-\beta_1^t},\quad
\widehat v_t=\frac{v_t}{1-\beta_2^t},\quad
\theta_t=\theta_{t-1}
-\eta\frac{\widehat m_t}{\sqrt{\widehat v_t}+\epsilon}.
\]
Here $t$ counts parameter updates, not microbatches. Calling \texttt{step()} after
each microbatch changes the parameters used for later forwards and changes the
moment histories. It no longer implements the fixed-$\theta$ derivation.

\Plate{05-adam}{The update boundary consumes the complete gradient once, advances
both moments and bias-correction clock, and writes the new parameters.}{evod27:adam}
\Listing{adam_step}
The oracle is out of place and validates tensor shape/dtype/device, finite values,
nonnegative second moments, and hyperparameter domains. It checks finite outputs
as well. Tests compare full-batch and accumulated PyTorch Adam updates,
every first/second moment, and the formula above. A two-step scalar calculation
independently checks that the bias-correction clock is not stuck at one.

Gradient parity alone can miss an optimizer bug. A meaningful equivalence gate is
$(L,g,m,v,\theta)$ parity under the same initial optimizer state. For restart tests,
the stored update count and moments are therefore part of the checkpoint, not
optional metadata.

\section{Clipping and loss scaling have an order}
If gradients are scaled by factor $s$, first accumulate the consistently scaled
contributions, then unscale the complete gradient, check finiteness, clip if
required, and perform one optimizer update. The update is skipped under an
overflow policy when nonfinite gradients are detected.

Clipping is nonlinear: generally
$\operatorname{clip}(g_1)+\operatorname{clip}(g_2)
\ne\operatorname{clip}(g_1+g_2)$.
For scalar threshold one, $g_1=2,g_2=-1.5$ gives zero after separate clipping,
but $0.5$ after summing and clipping. Microbatch-local clipping changes the
algorithm. Equally, clipping a scaled gradient at the unscaled threshold
quietly changes the threshold by $1/s$.

This chapter specifies this order but does not execute an accelerator AMP training
run. A production implementation must test actual scaler behavior, skipped-update
schedules, dtype policies, and optimizer states under its installed framework.

\section{Checkpointing changes retention, not the derivative}
Reverse-mode differentiation needs intermediate values to evaluate vector--Jacobian
products. \Term{Activation checkpointing} retains selected boundary inputs and recomputes
missing values during backward. It trades storage for extra execution; it must not
change the function being differentiated.

\Plate{03-checkpoint}{Forward retention and backward recomputation are different
phases of one derivative graph. Random state must replay the same stochastic
function, not merely produce tensors with matching shapes.}{evod27:checkpoint}
\Listing{checkpoint_grad}
The reference explicitly uses PyTorch's non-reentrant checkpoint path. For the
deterministic block, both loss and every parameter gradient match ordinary
execution in the Float64 fixture. The first-layer parameters still receive
gradients even though input data do not require gradients.

The PyTorch 2.10 documentation distinguishes the variants and describes default
recomputation checks based on shapes, dtypes, and devices \cite{torch-checkpoint}.
Those checks are not a numerical equality oracle. The tested runtime is
PyTorch 2.10.0; compiler modes and additional devices require separate validation.

\section{Same forward, wrong stochastic backward}
Dropout evaluates a function of both activations and a random mask. Backward
recomputation must reuse the appropriate stochastic state. With the same initial
CPU RNG state, ordinary execution, checkpointing with preservation, and
checkpointing without preservation have identical first forward losses.
Yet the last version recomputes with a different mask and produces a maximum
gradient discrepancy of about $0.220225$.

The preserved version matches the ordinary gradients and leaves the caller's CPU
RNG trajectory in the same post-operation state. This tests both the derivative
and the surrounding random-state contract. Model initialization uses a local
\texttt{fork\_rng} context and data use a dedicated generator; importing the
reference does not change the global default dtype or seed.

Random replay is only one hidden-state problem. A block that appends logs, updates
running statistics, writes an external memory, or mutates its input can behave
differently when executed again. Functionalization is preferable:
\[
(h_{\rm out},s_{\rm out})=F(h_{\rm in},s_{\rm in};\theta),
\]
with state transitions explicit and externally committed at the intended boundary.
A snapshot/restore strategy must account for every relevant state, not only the
visible tensor.

\section{Derive a checkpoint-placement toy before measuring a system}
Consider a chain of $L$ equal-size activation stages and blocks of at most $k$
layers. An illustrative inventory retains $\lceil L/k\rceil$ boundary slots and
at most $\min(k,L)$ slots while recomputing one block:
\[
I(L,k)=\lceil L/k\rceil+\min(k,L).
\]
Ignoring ceilings for $1\le k\le L$, $I\approx L/k+k$, whose derivative vanishes
at $k=\sqrt L$. This explains why neither one huge block nor tiny blocks necessarily
minimizes the toy inventory.

\Listing{chain_checkpoint_inventory}
\begin{center}\input{\Generated/chain.tex}\end{center}
For $L=64$, block size eight gives 16 toy slots; sizes four and sixteen both give
20. These slots are a declared upper inventory, not measured PyTorch saved tensors
or device bytes. At $k=1$ the toy reports 65 slots, so it must not be marketed as
an automatic improvement over retaining 64 ordinary stage outputs.

Checkpoint tradeoffs have a substantial research history \cite{chen2016}.
Real architectures have unequal tensor sizes, residual branches, recomputation
dependencies, attention matrices, and fused kernels. Our square-root derivation
does not determine their optimal placement or runtime speed.

\section{Mixed precision needs an operator-level policy}
\Plate{04-precision}{Low-precision compute, higher-precision reductions and optimizer
state, scaling, and finite-gradient gates form a policy. Simply casting every
tensor is not an equivalent training algorithm.}{evod27:precision}
FP16 has a smaller exponent range than FP32; BF16 retains an FP32-sized exponent
field but fewer significand bits. Loss scaling can help represent small gradients
in FP16; it cannot fix an incorrectly normalized objective or recover information
already lost in unstable arithmetic. Master weights and scaling are established
mixed-precision techniques \cite{micikevicius2018}, not a guarantee for every
optimizer/operator combination.

Sensitive reductions, normalization, attention logits, recurrence coefficients,
and optimizer moments deserve explicit dtype choices and reference comparisons.
Report finite-value incidence, skipped updates, loss/metric distributions across
seeds, and restart equivalence. A successful single training run is weak evidence
about numerical robustness.

\section{Apply the contracts to both model families}
DOGMA is the non-Transformer line. Its complete recurrent state, locus reads,
gates, and writes must preserve Chapter 22's ordering. Sequence chunking carries
state through time; microbatching partitions examples. Detaching a carried state
changes gradients and is truncated backpropagation, not activation checkpointing.

Hermon DNA is the Transformer line. Checkpointed blocks must preserve attention
masks, positional conventions, and stochastic policy. Decode-time KV cache is an
inference resource; ordinary training activation accounting should not quietly
substitute it for the actual saved graph.

Chapter 26's affine scan applies only to the stated input-conditioned recurrence.
Checkpointing that kernel does not establish that a state-dependent full DOGMA
controller can be parallelized identically. Evolutor must record which path was
selected and which state/numerical contract it actually implements.

\section{A measurement protocol that can fail}
\Plate{06-gate}{Semantic, numerical, resource, and task-quality gates are distinct.
Passing a CPU derivative test does not report an accelerator memory saving or a
genomic capability.}{evod27:gate}
Specify workload shapes, valid-token counts, dtype policy, optimizer, software,
hardware, and accumulation window. Warm up the complete step, including optimizer
state initialization. Reset peak statistics at a declared point, include forward,
backward, recomputation and update, and synchronize device execution around timing.

Distinguish live allocated bytes, allocator-reserved bytes, and external device
usage. Report repeat measurements and workload-dependent failure boundaries.
Selective retention, offload, and distributed sharding may move a peak rather than
remove it. A bandwidth bottleneck or extra communication can outweigh the saved
activation storage.

\section{Worked problems and implementation exercises}
\begin{enumerate}
\item Count the reference parameters. \textbf{Answer:} $128+16+48+3=195$.
\item Compute hypothetical FP32 Adam arrays. \textbf{Answer:} 3,120 bytes.
\item Compute all-Float64 arrays. \textbf{Answer:} 6,240 bytes, excluding overhead.
\item Why can two low-precision arrays still leave $16P$ bytes?
\textbf{Answer:} moments and an FP32 master copy can dominate the inventory.
\item Derive the global valid-weight mass. \textbf{Answer:} sum all fixed
nonnegative $w_{ij}$, including masks and importance weights.
\item Handle a zero-valid microbatch. \textbf{Answer:} its numerator contributes
zero divided by positive global $Z$.
\item Why reject target broadcasting? \textbf{Answer:} it can silently change
which output errors the objective includes.
\item Derive $\nabla_{W_2}L$. \textbf{Answer:} $D^\top H$, with
$D=2W\odot(P-Y)/Z$.
\item Scale all weights by seven. \textbf{Answer:} loss and gradient are unchanged.
\item Why is local-mean/$K$ wrong for masses 12,13,6?
\textbf{Answer:} it weights each microbatch equally rather than each weighted element.
\item When should Adam's clock advance? \textbf{Answer:} once per complete update.
\item What is missing from a gradient-only test? \textbf{Answer:} moment,
bias-correction, and parameter-update correctness.
\item Clip $2$ and $-1.5$ separately at one, then together.
\textbf{Answer:} separate sum zero; summed-and-clipped value $0.5$.
\item What precedes clipping under loss scaling? \textbf{Answer:} unscale and
check the complete accumulated gradient.
\item Why can equal checkpoint forward losses hide a bug?
\textbf{Answer:} backward recomputation can use a different random mask.
\item What does a shape/dtype/device check not prove? \textbf{Answer:} equal
recomputed values or equal derivatives.
\item Minimize $L/k+k$ for $L=64$. \textbf{Answer:} $k=8$, under the toy assumptions.
\item Is 16 toy slots measured peak memory? \textbf{Answer:} no.
\item Functionalize a DOGMA write. \textbf{Rubric:} explicit incoming state,
returned outgoing state, preserved read-before-write order, no duplicate external commit.
\item Design a deployment gate. \textbf{Rubric:} loss/gradient/moment/update parity,
RNG and restart contracts, realistic dtype failures, synchronized peak/time
measurement, repeated seeds, and explicit workload/hardware evidence.
\end{enumerate}

\section{Bridge to distributed ownership}
Chapter 28 assigns arrays and reductions to ranks. Global normalization,
optimizer-update boundaries, and stochastic-state provenance remain the same
questions, but now communication and ownership can invalidate them. An ideal
per-rank storage formula is only the beginning of that analysis.
